\emph{Preliminaries.} By Theorem \claimref{tower.ladder-pairs-satisfy-car-vectors} the modes $a_{x}^{*} = v_{x}$, $a_{x} = w_{x}$ ($x = (k, \alpha)$) satisfy $\{a_{x}, a_{y}^{*}\} = \delta_{xy}$, $\{a_{x}, a_{y}\} = 0$, and $v_{x}^{*} = w_{x}$. $J = g_{1}g_{2}$ has $J^{2} = -g_{1}^{2}g_{2}^{2} = -1$ and $J^{*} = (-g_{2})(-g_{1}) = g_{2}g_{1} = -J$; every plane generator anticommutes with $g_{1}$ and with $g_{2}$ (Theorem \claimref{std-algebra.orthogonal-anticommute}), so $J$ commutes with every mode, and $\mathrm{span}\{1, J\}$ is a field ($\cong \C$, conjugation $J \mapsto -J$) commuting with the mode algebra. Channel vectors and matrices have entries in it, with the positive Hermitian form $\langle u, u'\rangle = \sum_{\alpha}\bar u_{\alpha}u'_{\alpha}$; for a channel vector $u$ the element $c_{k}(u) = \sum_{\alpha}\bar u_{\alpha}a_{k,\alpha}$ is a mode, and orthonormal $u$'s give modes obeying the same relations. For a matrix $A$ on the one-particle index set (momenta of $I$ times channels) write $\hat A = \sum_{x,y}A_{xy}a_{x}^{*}a_{y}$. From the relations, $[a_{x}^{*}a_{y}, a_{z}^{*}a_{u}] = \delta_{yz}a_{x}^{*}a_{u} - \delta_{ux}a_{z}^{*}a_{y}$, hence
\[
[\hat A, \hat B] = \widehat{[A, B]}, \qquad (\hat A)^{*} = \widehat{A^{*}} ,
\]
the second because $(c\,a_{x}^{*}a_{y})^{*} = \bar c\,a_{y}^{*}a_{x}$. Now $\rho_{q}(M) = \widehat{S_{q}\otimes M}$, where $S_{q}$ is the partial shift with $(S_{q})_{k,k+q} = 1$ for $k, k + q \in I$ and zero otherwise, and $\rho_{q}(M')^{*} = \widehat{S_{q}^{T}\otimes M'^{*}}$.

\emph{(i).} $[S_{q}\otimes M, S_{q}^{T}\otimes M'^{*}] = S_{q}S_{q}^{T}\otimes MM'^{*} - S_{q}^{T}S_{q}\otimes M'^{*}M$, and $S_{q}S_{q}^{T}$, $S_{q}^{T}S_{q}$ are the diagonal projections onto $\{k : k + q \in I\}$ and $\{k : k - q \in I\}$. This is the window identity. Write the first set as $I$ minus the top $q$ momenta and the second as $I$ minus the bottom $q$:
\[
[\rho_{q}(M), \rho_{q}(M')^{*}] = \rho_{0}([M, M'^{*}]) - \sum_{k\ \mathrm{top}}a_{k}^{*}MM'^{*}a_{k} + \sum_{k\ \mathrm{bottom}}a_{k}^{*}M'^{*}Ma_{k} .
\]
On an edge-clean vector every top mode is empty, so $a_{k,\beta}$ annihilates it for $k$ in the top $q$; every bottom mode is filled, so $a_{k,\alpha}^{*}a_{k,\beta}$ acts on it as $\delta_{\alpha\beta}$. The last sum therefore acts as $q\operatorname{tr}(M'^{*}M)$. For $M = M'$ and a single channel this is the edge identity of Theorem \claimref{hawking.current-commutator-exactly-edge-split}.

\emph{(ii).} If $[M, M'^{*}] = 0$ the value on edge-clean vectors is the scalar $q\operatorname{tr}(M'^{*}M)$. Conversely let $X = [M, M'^{*}] \ne 0$; $\operatorname{tr}X = 0$. Since $N > q$ there is a bulk momentum $k_{0}$ (neither top nor bottom). Let $\psi_{0}$ be the edge-clean vector in which every momentum other than $k_{0}$ is filled in every channel or empty in every channel, and $k_{0}$ is empty, and consider $c_{k_{0}}(u)^{*}\psi_{0}$, one particle at $k_{0}$ in the channel state $u$. At a filled momentum $a_{k}^{*}Xa_{k}$ acts as $\operatorname{tr}X = 0$, at an empty one as $0$, and at $k_{0}$ it maps $c_{k_{0}}(u)^{*}\psi_{0}$ to $c_{k_{0}}(Xu)^{*}\psi_{0}$. If $\rho_{0}(X)$ were a scalar on all edge-clean vectors, every $u$ would be an eigenvector of $X$ with one eigenvalue, so $X$ would be a multiple of the identity matrix, and $\operatorname{tr}X = 0$ makes it zero. So $[\rho_{q}(M), \rho_{q}(M')^{*}]$ is a scalar on the edge-clean vectors if and only if $[M, M'^{*}] = 0$; for $M = M'$, if and only if $M$ is normal.

\emph{(iii).} The rotation $a_{k,\alpha} \mapsto e^{-\theta w_{\alpha}J}a_{k,\alpha}$ preserves the relations (the factors are unit elements of $\mathrm{span}\{1, J\}$ commuting with the modes), and $(e^{-\theta w_{\alpha}J}a_{k,\alpha})^{*} = a_{k,\alpha}^{*}e^{\theta w_{\alpha}J}$, so $a_{k,\alpha}^{*}a_{k+q,\beta} \mapsto e^{\theta(w_{\alpha} - w_{\beta})J}a_{k,\alpha}^{*}a_{k+q,\beta}$ and $\rho_{q}(M) \mapsto \rho_{q}(e^{\theta JW}Me^{-\theta JW})$. This equals $e^{h\theta J}\rho_{q}(M)$ for every $\theta$ exactly when every nonzero entry has $w_{\alpha} - w_{\beta} = h$, that is $[W, M] = hM$. Then a nonzero entry of $M^{n}$ needs $w_{\alpha} - w_{\beta} = nh$, impossible once $n|h|$ exceeds the spread of the weights: $M$ is nilpotent. A normal nilpotent matrix is zero: normality gives $|Mx|^{2} = \langle x, M^{*}Mx\rangle = \langle x, MM^{*}x\rangle = |M^{*}x|^{2}$, so $\ker M = \ker M^{*} = (\operatorname{im}M)^{\perp}$; if $M^{2}x = 0$ then $Mx \in \ker M \cap \operatorname{im}M = 0$, so $\ker M^{2} = \ker M$, and by induction $\ker M^{n} = \ker M$, which is everything. Hence for $h \ne 0$ and $M \ne 0$, $[M, M^{*}] \ne 0$ and by (ii) $\rho_{q}(M)$ is not a canonical boson on the edge-clean vectors; its commutator there is $q\operatorname{tr}(M^{*}M) + \rho_{0}([M, M^{*}])$. For the two chirality channels of a Dirac species the weights are $\pm\tfrac12$ (Theorem \claimref{synthesis.light-ray-good-part-of-a-spinor-is-two-chirality-arrays}(iv)), so the helicity-one matrices are the multiples of $E_{+-}$ and the helicity-minus-one matrices those of $E_{-+}$; for $M = E_{+-}$, $[M, M^{*}] = E_{++} - E_{--}$ and $\rho_{0}([M, M^{*}])$ is the difference of the numbers in the two channels.
For the doublet, the rotation acts as $\rho_{q}(M_{1}) \mapsto \cos\theta\,\rho_{q}(M_{1}) + \sin\theta\,\rho_{q}(M_{2})$, $\rho_{q}(M_{2}) \mapsto -\sin\theta\,\rho_{q}(M_{1}) + \cos\theta\,\rho_{q}(M_{2})$ exactly when $J[W, M_{1}] = M_{2}$ and $J[W, M_{2}] = -M_{1}$ (differentiate at $\theta = 0$; conversely the linear system integrates to the rotation). Put $M_{\pm} = M_{1} \mp JM_{2}$. Then $J[W, M_{\pm}] = M_{2} \pm JM_{1} = \pm J M_{\pm}$, so $[W, M_{\pm}] = \pm M_{\pm}$, and $M_{\pm}$ is nilpotent. If $[M_{i}, M_{j}^{*}] = 0$ for $i, j \in \{1, 2\}$ then $[M_{\pm}, M_{\pm}^{*}] = [M_{1} \mp JM_{2}, M_{1}^{*} \pm JM_{2}^{*}]$ is a combination of the $[M_{i}, M_{j}^{*}]$ and vanishes, so $M_{\pm}$ is normal and nilpotent, hence zero, and $M_{1} = \tfrac12(M_{+} + M_{-}) = 0$, $M_{2} = \tfrac{J}{2}(M_{+} - M_{-}) = 0$.
A matrix with $[W, M] = 0$ that is normal gives, by (ii), the scalar $q\operatorname{tr}(M^{*}M)$ on edge-clean vectors; the channel densities $E_{\alpha\alpha}$ are such.

\emph{(iv).} Let $M^{2} = 0$, $r = \operatorname{rank}M$, and $N \ge q$. Write $M^{*} = \sum_{s=1}^{r}\sigma_{s}t_{s}u_{s}^{*}$ with $\sigma_{s} > 0$ and orthonormal families $(u_{s})$, $(t_{s})$ (the singular value decomposition over $\mathrm{span}\{1, J\}$); $M^{*2} = 0$ gives $\sigma_{s}\sigma_{r'}\langle u_{s}, t_{r'}\rangle = 0$, so every $u_{s}$ is orthogonal to every $t_{r'}$. Then $B = \rho_{q}(M)^{*} = \sum_{k}\sum_{s}\sigma_{s}\,c_{k+q}(t_{s})^{*}c_{k}(u_{s})$, the sum over $k, k + q \in I$. Apply $B$ repeatedly to $\Omega$. A term with $k + q < 0$ vanishes on every vector reached, since its target mode at $k + q$ is filled in every channel and only modes at momenta $\ge 0$ change; a term with $k \ge 0$ vanishes, since the only filled modes at momenta $\ge 0$ lie in the span of the $t$'s, orthogonal to $u_{s}$. The remaining terms have $k \in \{-q, \dots, -1\}$ and define $\sigma_{k,s}^{+} = c_{k+q}(t_{s})^{*}c_{k}(u_{s})$ on $qr$ distinct sites $(k, s)$: they commute pairwise (even elements of distinct modes) and square to zero, and the vectors $\prod_{(k,s)\in S}\sigma_{k,s}^{+}\Omega$ for distinct sets $S$ are orthonormal. Hence $B^{p}\Omega = p!\sum_{|S| = p}\prod_{(k,s)\in S}\sigma_{s}\sigma_{k,s}^{+}\Omega$ and
\[
|B^{p}\Omega|^{2} = (p!)^{2}e_{p}(x), \qquad x = (\sigma_{s}^{2})_{(k,s)}, \qquad |B\Omega|^{2} = e_{1}(x) = q\operatorname{tr}(M^{*}M),
\]
with $e_{p}$ the elementary symmetric polynomial: zero for $p > qr$. Expanding $e_{1}(x)^{p}$, the products with distinct indices give $p!\,e_{p}(x)$ and for $p \ge 2$ the products $x_{i}^{p} > 0$ are extra, so $|B^{p}\Omega|^{2} < p!\,|B\Omega|^{2p}$ for $p \ge 2$, whereas a canonical boson has equality. For $M = E_{+-}$, $r = 1$, $x = (1, \dots, 1)$ ($q$ entries), $|B^{p}\Omega|^{2} = (p!)^{2}\binom{q}{p}$. On the same vectors the terms of $\rho_{q}(M) = B^{*}$ outside the sites vanish by the same occupancy argument, and $[\sigma_{k}^{-}, \sigma_{k}^{+}] = 1 - 2\sigma_{k}^{+}\sigma_{k}^{-}$ on the two states of a site, so $[\rho_{q}(M), B]B^{p}\Omega = (q - 2p)B^{p}\Omega$: the normalised commutator is $1 - 2p/q$, the quasi-spin of $q$ sites of Proposition \claimref{field-theory.how-close-tower-comes-ccr} (there with $n$ pairs). For the same-channel density at $q = 1$, $D = \rho_{1}(E_{++})^{*}$, the moves $k \to k + 1$ chain, and $|D^{2}\Omega|^{2} = 2 = 2!\,|D\Omega|^{4}$ (computed on the window $N = 2$).

\emph{(v).} $\omega([X, X^{*}]) = \omega(XX^{*}) - \omega(X^{*}X) = \omega(XX^{*}) \ge 0$ when $\omega(X^{*}X) = 0$, by positivity of $\omega$. For $q \ge 1$, $\rho_{q}(M)\Omega = 0$: each term removes a particle at $k + q$ and adds one at $k < k + q$, impossible in the sea since $k + q < 0$ forces $k < 0$, filled. By (i) on the edge-clean vector $\Omega$ (here $N \ge q$), $\omega_{\Omega}([\rho_{q}(M), \rho_{q}(M)^{*}]) = q\operatorname{tr}(M^{*}M) + \omega_{\Omega}(\rho_{0}([M, M^{*}]))$, and $\omega_{\Omega}(\rho_{0}(X)) = N\operatorname{tr}X = 0$; the reversed sea (filled for $k \ge 0$) is edge-clean for the reflected window and gives $-q\operatorname{tr}(M^{*}M)$, with $\rho_{q}(M)^{*}$ annihilating it. The timelike mode of the Gupta-Bleuler form in (K1) of Premise \claimref{synthesis.krein-form-premise-indefinite-inner-product-for-gauge-fields} annihilates the vacuum and has $[a_{0}, a_{0}^{\dagger}] = -1$, which the first sentence excludes for every pair mode and every state of the tower.

\emph{Checks.} The exact check verifies (i) on the window $N = 2$ with two channels for $q = 1, 2$ and nine pairs of six matrices (real and with $J$ entries, normal and nilpotent); (ii) at $q = 1$ on the edge-clean projector, including the values $q \pm 1$ of the helicity-one mode on two edge-clean states; (iii) the weights $\pm\tfrac12$, the doublet $M_{1} = E_{+-} + E_{-+}$, $M_{2} = J(E_{+-} - E_{-+})$ with $[M_{1}, M_{2}^{*}] \ne 0$, five integer helicity-one matrices on four channels, and the rotation at the half angle with $(\cos, \sin) = (\tfrac{12}{13}, \tfrac{5}{13})$; (iv) $q = 1, 2$, the Cartan comparison, and a rank-one nilpotent rational $M$; (v) $q = 1, 2$ on the sea and the reversed sea.
